\subsection{Stakeholder-Interview}\label{sec:stakeholder_interview}

Zur Ermittlung der unternehmerischen Anforderungen wurde zu Projektbeginn ein halbstrukturiertes Interview mit dem CEO des Unternehmens auf Basis eines Microsoft-Forms-Fragebogens (\hyperref[sec:anhang_a]{Anhang A}) geführt. Der Fragebogen umfasste acht thematische Fragen zu Zielen, Änderungswünschen, Pain-Points, Use-Cases, Priorisierungen, Funktionswünschen, KPIs und Sicherheit. Die Pain-Point-Matrix weist \emph{Nutzerabbrüche bei komplexen Prozessen}, \emph{mobile Nutzungshürden} und \emph{lange Ladezeiten} als kritischste Probleme aus (jeweils 5/6). Die Priorisierungsmatrix stuft \emph{Bessere Navigation}, \emph{Intuitivere Bedienung}, \emph{Modernes Erscheinungsbild} und \emph{Barrierefreiheit} als Höchstprioritäten ein (jeweils 6/6).

Aus den Ergebnissen leiten sich vier Handlungsfelder ab:
\begin{itemize}
    \item \textbf{Modernisierung des UI-Designs} -- Wunsch nach \emph{Angleichung an Standard-Windows-Verhalten} (Maus, Tastatur)
    \item \textbf{Verbesserung der Informationsarchitektur und Navigation} -- Pain-Point \enquote{Bessere Navigation} (6/6); Umsetzung des \emph{Top-Down-Prinzips} durch das in \autoref{sec:informationsarchitektur} entwickelte Navigationsmodell
    \item \textbf{Senkung der Support-Kosten} -- Reduktion der Pain-Points bei Nutzerabbrüchen (5/6), mobilen Hürden (5/6) und Support-Aufwand (4/6)
    \item \textbf{Stärkung der operativen Sicherheits-Architektur} -- \emph{Angriffsflächenminimierung, verteilte Schlüsselverwaltung und Security by Design} als strategische Vorgaben
\end{itemize}

Regulatorische Anforderungen sind als Zielvorgaben für Folgeiterationen definiert (vgl.\ \autoref{subsec:zielsetzung}).

Die beeinträchtigten Use-Cases umfassen die \emph{mobile Nutzung} (wird soweit möglich vermieden), die \emph{Rechtevergabe} (komplex, Gegenkontrollen erforderlich) und die fehlende Anpassbarkeit der Listenansicht. Der Erfolg des Redesigns wird an reduzierten Support-Anfragen, verringertem Schulungsaufwand und gesteigerter Neukundengewinnung gemessen.
